% Generated draft: requires mathematical and source review.
\documentclass[11pt]{article}
\usepackage[margin=25mm]{geometry}
\usepackage{fontspec}
\setmainfont{Latin Modern Roman}
\usepackage{amsmath,amssymb,mathtools,unicode-math}
\setmathfont{Latin Modern Math}
\usepackage{enumitem,needspace,xurl,hyperref}
\setlength{\parindent}{0pt}
\setlength{\parskip}{5pt}
\setlength{\emergencystretch}{3em}
\newcommand{\N}{\mathbb{N}}
\newcommand{\Z}{\mathbb{Z}}
\newcommand{\Q}{\mathbb{Q}}
\newcommand{\R}{\mathbb{R}}
\newcommand{\C}{\mathbb{C}}
\newcommand{\F}{\mathbb{F}}
\newcommand{\A}{\mathbb{A}}
\newcommand{\PP}{\mathbb{P}}
\newcommand{\E}{\mathbb{E}}
\newcommand{\eps}{\varepsilon}
\newcommand{\dd}{\,\mathrm{d}}
\begin{document}
\section*{UnsolvedMath ID: 800586}
\textbf{Source problem code:} catalogue entry
\subsection{Definitions and mathematical statement}
Let \(X\) be a topological space. An \(\omega\)-cover of \(X\) is a family \(\mathcal U\) of subsets of \(X\) such that \(X\notin\mathcal U\) and every finite subset of \(X\) is contained in some member of \(\mathcal U\). Under the standard countable-cover convention, write \(\mathrm B_\Omega(X)\) for the family of countable \(\omega\)-covers whose members are Borel subsets of \(X\). Here the Borel subsets form the sigma-algebra generated by the open subsets of \(X\).

For families \(\mathcal A\) and \(\mathcal B\) of covers of a space, the selection principle \(\mathsf S_1(\mathcal A,\mathcal B)\) means that for every sequence \((\mathcal U_n)_{n\in\N}\) of members of \(\mathcal A\), there are choices \(U_n\in\mathcal U_n\) such that \(\{U_n:n\in\N\}\in\mathcal B\). The principle \(\mathsf S_{\mathrm{fin}}(\mathcal A,\mathcal B)\) means that for every such sequence there are finite subfamilies \(\mathcal F_n\subseteq\mathcal U_n\) such that \(\bigcup_{n\in\N}\mathcal F_n\in\mathcal B\). Set
\[
P_1(X):=\mathsf S_1(\mathrm B_\Omega(X),\mathrm B_\Omega(X)),\qquad
P_{\mathrm{fin}}(X):=\mathsf S_{\mathrm{fin}}(\mathrm B_\Omega(X),\mathrm B_\Omega(X)).
\]
For a positive integer \(m\), give \(X^m\) the product topology, and interpret its Borel covers relative to that topology.

\textbf{Question, relative to the intended class of spaces.} For each \(i\in\{1,\mathrm{fin}\}\), does the implication
\[
\forall X\in\mathcal C\;\forall m\in\Z_{\geq1},\qquad
P_i(X)\Longrightarrow P_i(X^m)
\]
hold, where \(\mathcal C\) denotes the intended ambient class of spaces? These are two separate preservation questions; an affirmative answer to either would answer the source's question whether any of the two classes is preserved by finite powers.

\textbf{Unresolved scope.} The supplied excerpt does not specify \(\mathcal C\) or explicitly define \(\mathrm B_\Omega\). The source title provides a context of special sets of reals, but does not by itself determine the exact hypotheses. The definitions above use the standard interpretation as countable Borel \(\omega\)-covers and make the preservation question precise relative to \(\mathcal C\); the intended domain and cover convention require confirmation from the source.

\textbf{Supplied background qualification.} Write \(\Omega\) for the corresponding family of countable open \(\omega\)-covers, and \(\Gamma\) for the countably infinite open covers in which each point belongs to all but finitely many members. The supplied record states, with reference label \textup{[coc2]}, that \(\mathsf S_1(\Omega,\Gamma)\), \(\mathsf S_1(\Omega,\Omega)\), and \(\mathsf S_{\mathrm{fin}}(\Omega,\Omega)\) are the only properties in the open case preserved under finite powers. This is retained as a source-reported contextual assertion: the excerpt does not specify the collection of properties over which its word ``only'' ranges.
\subsection{Short English statement}
For each of two Borel cover selection properties, ask whether every positive finite Cartesian power of a space with that property also has it. The intended class of spaces and the countable Borel cover convention need source confirmation. The registry labels the problem open; this draft does not independently verify that status.
\subsection{Sources}
\begin{itemize}
\item Local UnsolvedMath registry: permanent ID \(800586\), exact source problem code \texttt{catalogue entry}, title ``Selection Principles in Mathematics --- Problem 12 ([CBC]).'' The supplied record lists Tsaban as proposer, gives the year 2006, and labels the problem open. Its upstream revision is \texttt{d371ee917ac9647fbb06cd6fc529deec6f3dda82}.
\item Supplied source reference: Tsaban, \emph{Selection Principles and Special Sets of Reals: Open Problems} (2006), Problem 12, \url{https://arxiv.org/abs/math/0609601}. The record reports extraction by \texttt{arxiv\_latex} and an access date of 2026-10-16. These are supplied provenance details; this draft does not claim to have accessed or independently verified the source.
\item The supplied attribution/context marker is \textup{[CBC]}. The contextual assertion about open covers carries reference label \textup{[coc2]}. Neither label is expanded bibliographically in the supplied record, so no author identities or additional references are inferred.
\end{itemize}
\end{document}
